Let $G=\Sym(\Omega)$, $|\Omega|=n\geq9$, have a C-group marking
$\rho_0,\ldots,\rho_{n-3}$. In the permutation graph each transposition
factor of $\rho_i$ is an edge labelled $i$. Edges of one label form a
matching. By the transitive-parabolic exclusion, every
$G_i=\gen{\rho_j:j\ne i}$ is intransitive. For each $i$ choose one
$i$-edge joining distinct $G_i$-orbits. Such an edge exists, since
otherwise $\rho_i$ would preserve every $G_i$-orbit and $G$ would be
intransitive. Let $F$ consist of these $n-2$ selected edges.
The fracture-forest method is due to Fernandes, Leemans and Mixer;
the following proof makes its complete use explicit.

\section{The possible extra edges}
\begin{lemma}\label{lem:forest}
The selected graph $F$ is a forest with two components. Every extra
edge joins its two components. Two extra edges intersect; if they have
different endpoint sets, they and two selected edges form an alternating
square. In that case they are the only extra edges.
Consequently the permutation graph is of one of three forms:
\begin{enumerate}
\item a tree with exactly one edge label occurring twice;
\item a tree with a nonempty set of parallel extra edges at one place,
      every repeated label occurring exactly twice;
\item a single alternating square, with all other edges forming attached
      trees and having pairwise distinct labels absent from the square.
\end{enumerate}
In the second case there are at least two extra edges; in all cases
every generator is a transposition or a product of two transpositions.
\end{lemma}
\begin{proof}
A cycle in $F$ would express the endpoints of any selected $i$-edge
as joined by a path using no $i$-edge, contradicting their distinct
$G_i$-orbits. Thus $F$ is a forest, and its $n-2$ edges give two
components. If an extra $i$-edge had both endpoints in one component,
its selected-tree path would have to contain no edge labelled $j\ne i$:
otherwise the remaining cycle would join the selected $j$-edge's
endpoints using labels different from $j$. But the only path consisting
solely of the single selected $i$-edge is that edge itself, which is
not an extra edge. Hence every extra edge crosses between components.

Take two extra edges $e_i,e_j$. If they are disjoint, their endpoints
are joined within the two components by nonempty selected paths. The
cycle formed by these paths and the extra edges can use no selected
label outside $\{i,j\}$. If $i=j$, the two nonempty paths already
require more than the one selected $i$-edge. If $i\ne j$, each path
has length at least two: a one-edge path between an endpoint of $e_i$
and an endpoint of $e_j$ can have neither label, since edges of a label
form a matching. Again there are too many selected edges of the allowed
labels. Thus the extra edges intersect, and in particular there cannot
be two extra edges of one label.

If $e_i=\{a,b\}$ and $e_j=\{a,d\}$ with $b\ne d$, the selected
path from $b$ to $d$ uses only $i,j$. The matching condition forces
exactly two edges, labelled $j,i$ in that order. This is the asserted
square. There cannot be a third extra edge. All crossing edges are
bipartite between the two components, so pairwise intersecting ones
have a common endpoint. A third edge parallel to either existing one
would force the same selected path to have two different final labels.
A third with another endpoint would give three vertices at pairwise
distance two in a tree. Their paths share one central vertex, but the
edge at a fixed outer vertex would again have two required labels.

If there is only one extra edge, the resulting simple graph is a tree.
If all extra edges are parallel, adjoining their common underlying edge
to $F$ also gives a tree. Otherwise the unique square just described is
the only cycle and its two labels occur exactly twice. This proves all
claims.
\end{proof}

\section{Elementary splitting identities}
All products below are ordinary permutations; disjoint factors commute.
\begin{lemma}\label{lem:splitting}
Let $\tau$ and $\sigma$ be disjoint transpositions.
\begin{enumerate}
\item If a transposition $e$ meets exactly one of their supports in one
point, then $\gen{\tau\sigma,e}$ contains both $\tau$ and $\sigma$.
\item For five distinct letters $a,b,u,v,w$, the three involutions
$\tau(vu),\tau(vw),(va)$, where $\tau=(ab)$, generate the full
symmetric group on those letters.
\item The subgroup $\gen{\tau(vu),(va)}$ does not contain $\tau$.
\end{enumerate}
\end{lemma}
\begin{proof}
For (1), the product of the meeting pair is a three-cycle, so cubing
$(\tau\sigma)e$ yields the other transposition; multiply back by
$\tau\sigma$ to obtain both. For (2), the product of the first two
involutions is a three-cycle on $u,v,w$. Conjugating $(va)$ by its
powers supplies $(ua)$ and $(wa)$. These generate the symmetric group
on $a,u,v,w$, including $(vu)$, and hence also $\tau$; adding $(ab)$
then gives the symmetric group on all five letters. For (3), both
listed generators preserve the unordered partition
$\{\{u,b\},\{v,a\}\}$, whereas $\tau$ does not.
\end{proof}

\begin{lemma}\label{lem:fivecores}
Suppose all repeated labels have the same transposition factor
$\tau=(45)$. Neither of the following marked five-letter configurations
can be extended by an edge from its support to a sixth letter while
preserving the intersection property:
\begin{equation}\label{eq:fivecores}
\begin{array}{c|ccc}
 &r_0&r_1&r_2\\\hline
\mathrm{path}&(12)\tau&(23)\tau&(34)\\
\mathrm{star}&(13)\tau&(23)\tau&(34).
\end{array}                                                    \end{equation}
The extending mark may be $(v6)$ for $1\leq v\leq5$, or
$\tau(v6)$ for $1\leq v\leq3$.
\end{lemma}
\begin{proof}
In the path configuration, $\tau\in\gen{r_0,r_2}$ by
Lemma~\ref{lem:splitting}(1). If the new mark $e=(v6)$ is untwisted,
then $\tau$ also belongs to $\gen{r_0,e}$ for
$v\in\{1,2,4,5\}$, or to $\gen{r_1,e}$ for $v=3$.
The common distinguished support is respectively $\{r_0\}$ or
empty, neither of which generates $\tau$. If $e=\tau(v6)$ with
$v=1$ or two, then $\tau\in\gen{r_2,e}$ by the same splitting
identity; the common support with $\{r_0,r_2\}$ is only $r_2$.
For $e=\tau(36)$, Lemma~\ref{lem:splitting}(2) puts $\tau$ in
$\gen{r_1,r_2,e}$, again with common support only $r_2$.

In the star configuration, Lemma~\ref{lem:splitting}(2) says
$\gen{r_0,r_1,r_2}=S_5$, which contains $\tau$. Every untwisted
$e=(v6)$ meets the support of at least one of $r_0,r_1$, and
Lemma~\ref{lem:splitting}(1) puts $\tau$ in that generator together
with $e$. The common support contains only that double transposition.
If $e=\tau(v6)$ with $v=1$ or two, then
$\tau\in\gen{r_2,e}$ as above. For $e=\tau(36)$,
Lemma~\ref{lem:splitting}(2) gives
$\tau\in\gen{r_0,r_2,e}$. The common support is
$\{r_0,r_2\}$, and Lemma~\ref{lem:splitting}(3) excludes $\tau$
from its subgroup. Every permitted extension therefore contradicts IP.
\end{proof}

\section{Eliminating parallel edges and long repetitions}
In the parallel or single-extra case of Lemma~\ref{lem:forest}, write
$T$ for the underlying tree, obtained by adding the common extra edge
$\tau=\{a,b\}$ to $F$. Each repeated mark is $\tau\sigma_i$;
all other marks are single transpositions.

\begin{lemma}\label{lem:distance}
The two transpositions in every repeated mark are separated in $T$ by
exactly one edge.
\end{lemma}
\begin{proof}
They are disjoint, so their distance is at least one. Suppose the path
between $\sigma_i$ and $\tau$ has at least two intervening edges.
Let $l$ be the first and $j$ the last. Since $j$ meets $\tau$, its
mark is a single transposition. If $l$ is also single, each of
$\gen{\rho_i,\rho_l}$ and $\gen{\rho_i,\rho_j}$ contains $\tau$
by splitting, whereas their common distinguished subgroup
$\gen{\rho_i}$ does not. If $l$ has repeated label, its mark is
$\tau\sigma_l$, with $\sigma_l=l$.

When there are at least three intervening edges, $j$ is disjoint from
both $\sigma_i$ and $\sigma_l$. Then each of
$\gen{\rho_i,\rho_j}$ and $\gen{\rho_l,\rho_j}$ contains the
symmetric group on the three letters supporting $\tau$ and $j$.
Their intersection is larger than $\gen{\rho_j}$, contrary to IP.
With exactly two intervening edges the local configuration is the path
in \eqref{eq:fivecores}, up to relabelling. Since $n\geq9>5$ and $T$ is connected,
there is an edge from this five-letter subtree to its complement.
Its label is new and has one of the eight forms in
Lemma~\ref{lem:fivecores}: an existing double label has already used
both of its matching edges, and the existing single label meets
$\tau$ and cannot be repeated. That lemma again contradicts IP.
\end{proof}

\begin{lemma}\label{lem:parallel}
There are no parallel edges in the permutation graph.
\end{lemma}
\begin{proof}
Suppose two distinct repeated marks are $\rho_i=\tau\sigma_i$ and
$\rho_j=\tau\sigma_j$. By Lemma~\ref{lem:distance}, each selected
edge $\sigma_i,\sigma_j$ has a single intervening connector to $\tau$.
These connectors are single transpositions, since they meet $\tau$.
If the connectors are distinct, the two selected edges are disjoint:
otherwise the underlying graph would have a cycle. Let the connectors
be $k,l$ for $\sigma_i,\sigma_j$, respectively. The connector $k$
meets $\tau$ and is disjoint from $\sigma_j$, while $l$ is disjoint
from $\sigma_i$. Splitting therefore puts $\tau$ in both
$\gen{\rho_j,\rho_k}$ and $\gen{\rho_i,\rho_l}$. Their distinguished
supports are disjoint, contradicting IP.

If the connectors coincide, the local configuration is the star in
\eqref{eq:fivecores}. Again connectedness and $n>5$ supply an edge to a sixth letter,
and Lemma~\ref{lem:fivecores} excludes every permitted label.
\end{proof}

\section{The complete core classification}
\begin{theorem}\label{thm:classification}
Every rank-$n-2$ C-group marking of $S_n$, $n\geq9$, is, up to
relabelling letters and generators, one of the three constructions:
\begin{enumerate}
\item core $(23),(12)(34)$, with a tree attached at the core vertex $4$;
\item the same core, with a tree attached at vertex $3$;
\item core $(12)(34),(13)(24)$, with a tree attached at vertex $4$.
\end{enumerate}
In each case the core vertex is a leaf of the attached tree, its first
edge is a single transposition, and every tree edge has a different
new label. There are no other cases.
\end{theorem}
\begin{proof}
By the preceding lemmas the graph is either a simple tree with one
repeated label at distance one, or a unique alternating square. The
first gives a four-letter core consisting of two disjoint transpositions
with a single connector; it is the first listed core, whose marked
group is the transitive dihedral group $D_8$. The square gives the
transitive Klein four core in the third construction.

For any transitive group $H\leq\Sym(Q)$ and a new letter $x$, the group
$\gen{H,(vx)}$ is $\Sym(Q\cup\{x\})$: conjugation by $H$ gives all
star transpositions $(qx)$, $q\in Q$. If two distinct edges leave
one of our four-letter cores, with respective new endpoints $x,y$,
the two groups generated by the core and one edge therefore intersect
in $\Sym(Q)$. IP would require their intersection to be the core group,
which is only $D_8$ or $V_4$. This is impossible. Hence exactly one edge
leaves the core, and all remaining edges form one attached tree.

The marked dihedral core has two vertex orbits under permutations
centralizing both of its distinguished generators: the two endpoints
of the length-three path, and its two middle vertices. They give
constructions A and B. To verify that these orbits are distinct, the
single core transposition fixes exactly the two endpoints, whereas
its support is the middle pair. Reversing the path interchanges the
two vertices in each pair while preserving both marks. In the square
case the Klein four group centralizes its own two marks and is
transitive on the four vertices, giving only construction C.
This proves exhaustion. The core-attachment theorem of [U] and its local verification above establish
the full intersection property of every listed construction, and the
flag-completion proof below establishes the regular-hypertope conclusion.
\end{proof}

\section*{Exact finite validation and source scope}
The program \texttt{local\_obstructions.py} independently enumerates
subgroups of $S_6$ for all sixteen extensions in \eqref{eq:fivecores} and the two
six-letter cases with distinct connectors. It emits explicit words
for the common element $\tau$ and checks its exclusion from the
subgroup on the common labels. Five deliberately corrupted certificates
are rejected. These are validation examples: the splitting identities
above prove all cases uniformly, so no finite enumeration is a premise
of Theorem~\ref{thm:classification}. Completeness of the local cases
is established by the forest and distance arguments, not inferred from
the program's case count.

The theorem recovers the necessity assertion of Fernandes--Leemans,
Theorem 1.2, in the 2018 accepted manuscript, together with the
intransitive-parabolic reduction proved separately. The fracture method
and the three historical families are theirs. The module and simple-kernel bounds,
explicit local splitting proof and decorated-coset flag completion are
the present proof organization; no historical conjecture solution or
priority claim follows from this recovery.
